\section{Introduction}\label{sec:intro}

kNN search asks a simple but fundamental question: given a query object, which $k$ objects in a reference collection are most similar to it? In modern data systems, objects are commonly represented as feature vectors, and similarity is measured by a distance function over those vectors. A kNN join lifts this operation from one query to a whole query set: for every object in a query set $U$, it maintains the $k$ closest objects in a reference set $I$~\cite{bohm2003supporting,bohm2004k,xia2004gorder,yu2007efficient}.

Existing dynamic kNN join studies mainly consider one-sided dynamic settings, where one input set is updated while the other side is fixed or plays a static role~\cite{yang2014continuous,ukey2022efficient,ukey2023efficient,ukey2023knn}. This paper studies a more demanding dynamic version of the problem, which we call dual-dynamic kNN join. In this setting, both sides of the join evolve: new query objects may arrive, old query objects may be removed, new reference objects may be inserted, and existing reference objects may disappear or be retired. The goal is to maintain the kNN join table as both $U$ and $I$ change. To the best of our knowledge, this is the first work to study kNN join maintenance under fully dynamic updates on both sides.

Dual-dynamic kNN join is useful whenever two evolving collections of vectorized entities must be compared continuously. In healthcare, a newly admitted patient or an updated patient record can be treated as a query object in $U$, represented by diagnoses, laboratory tests, medication history, imaging features, and temporal clinical indicators. The reference side $I$ may contain historical patient cohorts, treatment-response profiles, drug-safety cases, or genomic records, which also change as hospitals correct records and medical databases are refreshed. Maintaining the neighborhood relation between current patients and these evolving references can support case reanalysis, similarity-based risk prediction, and diagnosis assistance, where unstable or incorrect neighbor relations may affect downstream clinical decisions~\cite{an2023comprehensive,fang2021patient,gottlieb2013method}. In finance, current transactions, accounts, or market states can form the query side $U$, encoded by behavioral and temporal features, while confirmed fraud cases, legitimate comparison cases, risk patterns, or market observations form the reference side $I$. Analysts need to track which new activities are close to suspicious historical patterns, while the reference collection itself changes as new fraud cases, legitimate cases, and market observations are added~\cite{al2021financial,immadisetty2025real,li2023artificial,gao2024machine}. Similar needs also arise in recommendation, anomaly detection, security monitoring, scientific data analysis, and vector-database workloads, where embedding-based objects and their reference collections are continuously refreshed~\cite{ahuja2019movie,tian2014anomaly,li2007network,wang2020log,pan2024survey}.

The problem is difficult for three reasons.

\noindent(1) Contemporary data volumes and feature dimensions keep increasing. Rich representations encode many attributes of an entity, but each additional dimension makes distance computation more expensive. High dimensionality also weakens many coarse geometric filters because pairwise distances tend to concentrate~\cite{beyer1999nearest,kouiroukidis2011effects,gao2023high,zhao2023towards}.

Approximate kNN techniques are a common way to reduce high-dimensional search cost~\cite{bawa2005lsh,andoni2008near,jegou2010product,muja2014scalable,li2019approximate,gao2023high,zhao2023towards}. However, when a fully dynamic kNN join must be maintained exactly, approximate methods are not suitable because approximation changes the maintenance contract. An update may require both kNN computation and affected-query localization. If each approximate stage has a small probability $\epsilon$ of missing a true neighbor or an affected query, then even one query-side repair followed by one reference-side localization already creates multiple opportunities for the materialized table to deviate from the exact result; under a simple independent-error view, preserving correctness through two such stages has probability at most $(1-\epsilon)^2$, and the risk compounds across long update streams. This is especially problematic in high-stakes settings such as clinical decision support or financial monitoring, where a small missed neighborhood relation can trigger a much larger downstream error.

\noindent(2) The maintained result is fully dynamic at the view level. A single insertion or deletion can potentially change the kNN neighborhood states of many queries. Under exact high-dimensional maintenance, brute-force recomputation of the entire join after each update is not a viable option in update-intensive workloads. The setting therefore calls for a maintenance algorithm that can accurately locate the rows that may change and repair only those rows.

\noindent(3) Efficient kNN join processing usually relies on indexes~\cite{bohm2001cost,bohm2004k,xia2004gorder,yu2007efficient}. In a highly dynamic setting, these indexes must not only accelerate search but also support repeated updates. This already makes dynamic indexing a core difficulty. Moreover, updates on both sides of the join may require more complex, and possibly multiple, index structures to accelerate different search operations. The increased index complexity further deepens the dynamic-maintenance challenge, because the indexes must be kept synchronized with the materialized join table throughout the update stream.

To address these challenges, we propose \dualtree, a framework centered on a dual-index structure. A \deltatree over the reference set $I$ accelerates exact kNN search from queries in $U$ to references in $I$, while an \hdrtree over the query set $U$ accelerates RkNN search from an updated reference in $I$ back to the queries whose join rows may change. Together, these indexes provide the two high-dimensional search directions required by fully dynamic kNN join maintenance.

Supporting these searches throughout an update stream requires both index components themselves to remain current. \deltatree already supports updates to its indexed reference set, whereas the original \hdrtree assumes a static query set and does not provide the complete query-side update mechanisms required here. We extend \hdrtree with query insertion and deletion, together with explicit local maintenance of its cluster-level pruning parameters when changed query rows alter their kNN distances. This extension yields a dual index that can evolve with both $U$ and $I$.

Using this dynamically maintainable dual index, we formulate fully dynamic kNN join as the coupled maintenance of two evolving structures: the materialized join table and the dual-index state. We develop four strategies for query insertion, query deletion, reference insertion, and reference deletion. Each strategy invokes the appropriate kNN or RkNN search direction to compute or identify the rows that require change, and orders table and index updates according to their data dependencies. After every update, the join table remains exact and both index components remain current, allowing subsequent searches to continue returning correct results.

To further address the high-dimensional challenge, we introduce a fine-grained low-dimensional pruning stage at the leaf nodes of \dualtree. This stage prunes visited points before full-dimensional verification whenever the low-dimensional distance is already sufficient to rule them out. We further derive a cost-aware dimension-selection model that balances pruning power against projection cost, allowing the pruning dimensionality to be selected by an explicit optimization objective rather than by a heuristic setting.

In summary, this paper makes the following contributions:

\noindent(1) \emph{Bidirectional search through dual indexing.}
We propose \dualtree, which combines \deltatree-based exact kNN search over the reference set with \hdrtree-based RkNN search over the query set. This dual-index structure accelerates both search directions required by fully dynamic high-dimensional kNN join maintenance.

\noindent(2) \emph{Dynamic dual-structure maintenance.}
We extend \hdrtree with query-side structural updates and explicit local maintenance of its pruning parameters. On the resulting dynamically maintainable dual index, we formulate each fully dynamic update as coupled maintenance of the materialized join table and the dual-index state, and provide four dependency-aware strategies that keep both structures synchronized.

\noindent(3) \emph{Fine-grained pruning and cost-aware modeling.}
We extend reduced-space pruning from the cluster level to the point level at leaf nodes and derive a cost-aware model that selects the projection dimensionality by balancing pruning power against projection cost. This replaces heuristic dimension selection with an explicit optimization objective.

\noindent(4) \emph{Comprehensive evaluation.}
We evaluate the proposed methods on eight real high-dimensional datasets with dimensionalities from 128 to 4096. The results show that \dualtree consistently outperforms the non-Dual-Tree baselines, while \dualtreeopt achieves an average speedup of 1.88$\times$ and up to 2.65$\times$ over the fastest non-Dual-Tree baseline on each dataset.

The remainder of the paper is organized as follows. Section~\ref{sec:related-works} reviews related work. Section~\ref{background} gives definitions and index preliminaries. Section~\ref{sec:framework} presents the dual-index framework, the dynamic extension of \hdrtree, and the coupled maintenance strategies. Section~\ref{sec:pruning} develops fine-grained pruning and cost-aware dimension selection. Section~\ref{sec:exp} reports the experimental results, and Section~\ref{sec:conclusion} concludes the paper.
